\documentclass[11pt]{article}

\usepackage[margin=1in]{geometry}
\usepackage{amsmath, amssymb, amsfonts}
\usepackage{booktabs}
\usepackage{graphicx}
\usepackage{hyperref}
\usepackage[numbers,sort&compress]{natbib}

\title{Recoverable Urban Resilience: Learning the Laws of Managed Recovery from Empirical Urban Disruptions}
\author{Draft}
\date{\today}

\begin{document}
\maketitle

\begin{abstract}
Urban resilience is often measured as the observed ability of a city to withstand disruption and return toward normal function. This paper reframes resilience as a decision-sensitive property: how much of the observed functional loss is recoverable through constrained, timely, and equitable management intervention. We combine empirical urban disruption data, a data-calibrated functional recovery optimization model, and cross-city data mining to identify when managed recovery has high leverage.
\end{abstract}

\section{Introduction}
Urban recovery is usually analyzed through observed degradation and rebound. The central question is whether the system returned to normal, how quickly it did so, and which places experienced the largest or longest deficits. This paper asks a different question: what portion of the observed loss was actually recoverable through intelligent management under realistic constraints?

This framing separates observed resilience from recoverable resilience. A city can recover quickly because endogenous dynamics are strong, even when management decisions have little additional leverage. Another city can appear fragile in observed data while still containing highly recoverable losses if targeted intervention can reduce exposure or restore access.


\section{Conceptual Framework}
We define recoverable urban resilience as the fraction of cumulative functional loss that can be reduced by optimized intervention relative to a no-additional-intervention baseline. The concept links three layers: underlying functional deficit, locally experienced deficit, and access-weighted resident loss.

The data-mining stage evaluates whether empirical mobility, rainfall, demand, and network data contain the ingredients needed to estimate those layers before the counterfactual optimization model is introduced.


\section{Data}
The empirical foundation consists of multi-city U.S. mobility and network data. Current raw data cover 11 U.S. cities with demand and network outputs, 7 cities with usable large-scale speed observations, and 3 cities with existing preliminary resilience-index files. The raw data include speed observations, rainfall time series, TMC identification tables, demand matrices, road-network nodes and links, and traffic-assignment performance outputs.

Large raw files are stored outside git. Reproducible scripts generate compact summaries, figures, and tables for the manuscript.


\section{Functional Recovery Optimization}
The optimization model allocates high-level recovery primitives across space and time to reduce cumulative access-weighted functional loss. Durable restoration reduces underlying deficit, temporary capacity reduces locally experienced deficit, and substitution or control reduces origin-level experienced loss.

Let $b_{i,t}$ denote underlying functional deficit, $d_{i,t}$ local experienced deficit, and $\ell_{i,t}$ access-weighted functional loss. A functional dependence matrix $Q_t$ maps destination deficits to origin-level loss. The model minimizes $\sum_t \sum_i p_i \ell_{i,t}$ subject to recovery dynamics, resource budgets, response delays, and bounded functional states.


\section{Data Mining Evidence}
The first data-mining pass evaluates whether the available data can support the recoverable-resilience research program before the optimization model is introduced. The analysis covers data availability, rainfall shock characterization, speed-deficit and recovery proxies, spatial concentration of deficit burden, demand-network structure, and alignment between preliminary resilience metrics and network exposure.

Across the usable speed datasets, the strongest observed functional-loss signals appear in Chicago and New York. Chicago has a high p90 speed deficit and a large share of observations with at least 20 percent speed deficit, while New York has a weaker but still substantial deficit signal over a much larger observed TMC set. Rainfall-speed alignment is positive but moderate: San Antonio, Austin, New York, and Chicago show the clearest lagged association between rainfall and speed deficit, with event-level recovery proxies typically on the order of a few hours in the current monthly data.

The spatial concentration results are important for the paper's decision-centered framing. In several cities, the top 10 percent of TMCs account for roughly 30--40 percent of the total speed-deficit burden. This does not prove recoverability, but it shows that observed loss is not uniformly diffuse. A targeted intervention model may therefore have meaningful spatial leverage to test.

Demand and network data provide a baseline empirical foundation for functional dependence. All 11 cities have OD demand and link-performance outputs, with strong cross-city differences in OD scale, destination concentration, and congested-volume share. These data can support a baseline dependence matrix \(Q_t\), although destination service importance \(S_j\) still requires additional POI, employment, health-care, retail, population, or other functional-exposure data.

Finally, for the three cities with existing resilience-index files, the observed loss metrics can be matched to link-performance exposure. Chicago shows the clearest positive alignment between loss magnitude and link volume; New York shows weaker but positive alignment; Houston is weak or mixed. This suggests that the data contain functional-exposure signal, but descriptive alignment alone cannot answer how much loss would be recoverable under alternative intervention allocations.


\section{Discussion}
The discussion will distinguish what the data can identify directly from what requires counterfactual modeling. Descriptive data mining can reveal disruption, heterogeneity, endogenous recovery, and functional dependence. It cannot by itself identify the additional recovery achievable under unobserved interventions.


\section{Conclusion}
This draft supports the recoverable-resilience framing. The empirical data reveal rainfall-aligned mobility disruptions, nonuniform speed-deficit burden, OD functional dependence, and heterogeneous natural recovery. The optimization model then estimates the counterfactual portion of this loss that is recoverable under a standardized management regime.

The learning analysis extracts two linked laws from the optimizer-derived recovery-value field. First, local marginal recovery value is activated when persistent future deficit overlaps with OD exposure and feasible intervention efficiency. Second, finite-budget allocation must be residual: after each deployment, the remaining value depends on the current unrecovered loss, remaining budget, remaining time, and diminishing returns. At the event level, decision-criticality is governed by the top tail of marginal recovery values rather than disruption magnitude alone.

The current results should be read as regime-conditional evidence, not as a final universal law. The next strongest extensions are wider scenario-specific LP closure, intervention-parameter sensitivity, and a graph-based surrogate for larger-scale law extraction.


\bibliographystyle{plainnat}
\bibliography{references}

\end{document}
